\section{¿Los modelos preentrenados ya saben razonar?: de la greedy decoding a la CoT decoding}

Si las trayectorias intermedias son útiles, el siguiente paso no es hacer fine-tuning de inmediato, sino revisar primero si esas trayectorias ya existen en la distribución preentrenada. La clase plantea una hipótesis contraintuitiva: tal vez el modelo no sea «del todo incapaz», sino que la greedy decoding eligió demasiado pronto una rama corta, de alta probabilidad local, pero errónea. Así, el objeto de estudio pasa de la capacidad de los parámetros a la estrategia de decodificación.

\subsection{«No sabe» puede significar solo que el primer candidato es incorrecto}

Las diapositivas presentan primero una creencia común en ese momento: un LLM preentrenado necesita ingeniería de prompts o fine-tuning para producir razonamiento. La página siguiente la refuta directamente y subraya que «all we need is decoding». Al leer estas dos páginas hay que conservar los matices: esto no significa que cualquier modelo preentrenado pueda resolver cualquier tarea, sino que \textbf{no basta con que falle la salida voraz top-1 para afirmar que en su distribución no existe una trayectoria correcta}.

\lecturefigure{slide-06.jpg}{Creencia común: sin prompting ni fine-tuning, un modelo preentrenado no puede razonar}{6}

\lecturefigure{slide-07.jpg}{Contrapropuesta de la clase: la distribución preentrenada puede contener ya trayectorias de razonamiento}{7}

\teachervoice{00:06:20--00:08:30, el docente desplaza el énfasis de «volver a entrenar el modelo» a «revisar primero la distribución de generación». Las diapositivas advierten: latent capability y top-1 behavior no son el mismo concepto; una salida fallida puede deberse a la búsqueda y al ordenamiento, o a que en la distribución no existe en absoluto una solución correcta, y ambos casos deben distinguirse mediante el análisis de candidatos.}

\subsection{El problema de las manzanas: generación de candidatos, ordenamiento erróneo y confianza en la respuesta}

Las tres páginas siguientes descomponen un fallo en dos etapas: «si el candidato existe» y «si el ordenamiento es correcto». En el problema de las manzanas, la salida greedy es 5, pero en el conjunto de candidatos ya aparece la trayectoria correcta «el papá tiene 5, en total $3+5=8$». La página 8 demuestra primero que revisar más candidatos permite encontrar la solución correcta; la página 9 descarta después un atajo tentador, elegir la respuesta según la longitud de la respuesta generada, porque un texto largo también puede equivocarse con seguridad; la página 10 propone finalmente ordenar según \textbf{la confianza del modelo en la parte de la respuesta final}. Al leer las figuras hay que seguir cómo un mismo grupo de candidatos se reordena en las tres páginas con distintas reglas de selección, en lugar de tomar las tres páginas como capturas repetidas.

\lecturefigure{slide-08.jpg}{Cuando greedy falla, el conjunto de candidatos aún puede contener un razonamiento correcto}{8}

\lecturefigure{slide-09.jpg}{La longitud no es un indicador confiable de la calidad del razonamiento}{9}

\lecturefigure{slide-10.jpg}{Seleccionar candidatos según la confianza en la respuesta final, no según la longitud de la trayectoria}{10}

Para un candidato $y^{(k)}=(r^{(k)},a^{(k)})$, puede definirse la log-probabilidad promedio del intervalo de la respuesta:
\[
s_k=\frac{1}{|a^{(k)}|}\sum_{j=1}^{|a^{(k)}|}
\log p_\theta\!\left(a^{(k)}_j\mid x,r^{(k)},a^{(k)}_{<j}\right),
\qquad k^*=\arg\max_k s_k.
\]
Aquí $s_k$ es la confianza normalizada de la respuesta del candidato $k$, $|a^{(k)}|$ es el número de tokens de la respuesta y $k^*$ es la selección final. Considerar solo el último token es adecuado para respuestas tipo etiqueta; para números de varios tokens, programas o texto libre deben usarse límites explícitos de la respuesta y normalización por longitud; de lo contrario, las respuestas cortas obtienen una ventaja injusta.

\begin{warningbox}{La confianza no es prueba de corrección}
Las probabilidades del modelo provienen de su propia distribución y pueden tener un exceso de confianza sistemático; la extracción de la respuesta también puede cortar mal las unidades, los signos o el cuerpo de un programa. La confianza sirve como señal para ordenar candidatos, pero no debe cargar por sí sola con la detección de hallucination, los umbrales de puesta en producción ni las garantías de seguridad.
\end{warningbox}

\subsection{El algoritmo de dos pasos de la Chain-of-Thought Decoding}

La estructura de la CoT decoding es sencilla: primero, ir más allá de greedy y explorar varios candidatos generados; después, elegir según la confianza en la respuesta final. Es distinta del CoT prompting: la primera modifica \textbf{la decodificación y el ordenamiento}; el segundo modifica \textbf{el prompt de entrada}. Un mismo modelo puede aplicar CoT decoding sin ejemplos, o seguir usando greedy con un CoT prompt; los dos ejes deben mantenerse separados.

\lecturefigure{slide-11.jpg}{Chain-of-Thought Decoding: ampliar los candidatos y luego ordenarlos según la confianza en la respuesta}{11}

\begin{lstlisting}[language=Python,caption={Pseudocódigo mínimo para reordenar candidatos CoT según la confianza en la respuesta}]
def cot_decode(model, prompt, samples, answer_parser):
    candidates = model.sample(prompt, n=samples)
    ranked = []
    for response in candidates:
        answer_span = answer_parser(response.text)
        if answer_span is None:
            continue
        score = mean(response.logprobs[answer_span])
        ranked.append((score, response.text))
    if not ranked:
        raise ValueError("no parseable answer")
    return max(ranked)[1]
\end{lstlisting}

\begin{knowledgebox}{Lista de verificación de ingeniería: qué requiere ordenar candidatos}
Deben registrarse la temperatura de muestreo y el top-$p$, el número de candidatos, la versión del analizador de respuestas, la log-prob del intervalo de la respuesta, la proporción de candidatos repetidos y la tasa de análisis inválidos. De lo contrario, la «mejora por decodificación» podría deberse solo a gastar más tokens, a un cambio de formato o a una filtración accidental de los límites de la respuesta, y el experimento no podría reproducirse.
\end{knowledgebox}

\teachervoice{00:10:24--00:12:10, el docente enfatiza que la puntuación considera la probabilidad condicional de la final answer, y no se obtiene multiplicando las probabilidades de toda la trayectoria larga para compararlas directamente. Nota de clase: cuanto más larga es la trayectoria, menor es por naturaleza su probabilidad conjunta; si no se delimitan la respuesta ni se normaliza, el sesgo de longitud se confunde con calidad.}

\subsection{Resumen de la sección}

Esta sección descompone la afirmación «los modelos preentrenados no saben razonar» en dos proposiciones verificables: si la trayectoria correcta existe en la distribución de candidatos y si el decodificador puede colocarla en el primer lugar. La CoT decoding resuelve el segundo problema, pero depende del análisis de la respuesta y de la calibración de las probabilidades; si en la distribución falta por completo la trayectoria correcta, es necesario recurrir al prompting o al entrenamiento.
